\chapter{Literature Review}
\label{ch:ch2}

This chapter reviews the literature that informs and positions the present study. The review is organised around six thematic streams: the segmentation of elderly populations using survey data (Section 2.1), the digital divide in ageing societies (Section 2.2), the distinction between objective conditions and subjective well-being in later life (Section 2.3), the relationship between cognitive functioning and digital engagement (Section 2.4), the comparative welfare regime framework relevant to the cross-country component (Section 2.5), and the under-used SHARE Expectations module (Section 2.6). Section 2.7 synthesises the identified gaps and articulates the positioning of this thesis relative to the existing body of work.

\section{Segmentation of Elderly Populations}
The heterogeneity of the elderly population has long been recognised as a challenge for health and social policy. Older adults differ not only in their clinical conditions but in their functional capacity, social resources, economic circumstances, and subjective experiences. Population segmentation---the identification of subgroups with internally homogeneous and externally distinct profiles---offers a systematic approach to mapping this heterogeneity and informing targeted interventions.

\emph{\textbf{2.1.1 Needs-based segmentation in healthcare}}

The most established stream of elderly segmentation originates from healthcare resource allocation. Needs-based population segmentation models classify individuals into categories defined by their expected healthcare utilisation, enabling policymakers to allocate resources proportionally to the complexity of care needs (Lynn et al., 2007). The Bridges to Health model, for instance, proposes eight segments ranging from `healthy' to `frail' to `dying', each associated with distinct care goals and service configurations. More recently, Vuik et al. (2016) developed and applied a cross-country segmentation tool (CCSST) to SHARE data, classifying over 270,000 observations from 14 European countries into five global impression segments based on clinical indicators, self-reported diagnoses, and performance-based tests. Their analysis demonstrated significant cross-country variation in the distribution of need segments, with Italy and Israel exhibiting the highest proportions of high-need individuals.

Vuik et al.'s (2016) approach deserves closer scrutiny. Their five segments were defined by expected resource consumption rather than by the multidimensional profile of the individual---a distinction that matters for the policy uses to which a segmentation is subsequently put. A resource-consumption segmentation answers the question 'how much will this group cost?' while the segmentation developed in this thesis answers the logically prior question 'who are they and what do they need?', and the two framings are not reducible to one another. Lafortune et al. (2009), in an OECD working paper that became an influential reference for the field, argued that needs-based segmentation is most valuable when it integrates functional, cognitive and social dimensions; yet the subsequent empirical literature has largely retreated to a clinical-resource framing, a gap the present thesis attempts to close.

The World Health Organization's (2015) World Report on Ageing and Health provided the most influential high-level articulation of this broader view, proposing 'intrinsic capacity'---the composite of physical, mental and functional capacities that individuals can draw upon---as the organising concept for elderly-oriented policy. The WHO framework is conceptually aligned with the multidimensional segmentation approach adopted here, but it was proposed as a policy vocabulary rather than operationalised at the individual level; operationalising it at the individual level requires precisely the kind of profile analysis this thesis undertakes. Bock et al. (2016), reviewing clustering applications in geriatric populations, noted that the standard clinical segmentations tend to produce solutions highly sensitive to the inclusion or exclusion of single indicators---a fragility that contrasts with the cross-method and cross-country stability documented in Chapter 6 of the present work.

More substantively, the limits of a purely clinical segmentation become visible when one considers the counterfactual profile structure such an approach would produce. Without subjective and social variables, the Moderate Isolated profile---28.5\% of the Italian sample and the single most policy-relevant segment of this thesis---would not emerge as a distinct cluster. A clinical-only segmentation would classify these individuals as belonging to a generic 'low-need' group, an inversion of their actual vulnerability that illustrates why the choice of input dimensions is not a technical preference but a substantive determinant of what the segmentation is able to see. The empirical demonstration of this point is provided in Chapter 6, where a four-dimensional specification is compared head-to-head with the five-dimensional solution and is shown to halve the external discriminative power of the clustering.

While these approaches are valuable for healthcare planning, they share a fundamental limitation: they segment populations exclusively on the basis of clinical and functional indicators. Social connectivity, cognitive functioning, digital accessibility, and subjective well-being---dimensions that are increasingly recognised as determinants of health trajectories---are absent from the segmentation framework. The resulting profiles, while clinically coherent, provide an incomplete picture of the elderly person's actual situation and needs.

\emph{\textbf{2.1.2 Multidimensional segmentation on SHARE}}

A growing body of literature has employed SHARE data for multidimensional segmentation of the elderly, extending beyond purely clinical indicators. Lestari et al. (2022) applied Latent Class Analysis to SHARE data, identifying four classes based on health status and economic resources. Their work demonstrated the feasibility of person-centred segmentation on SHARE's rich variable set, but limited the analysis to two dimensions, excluding social, cognitive, and subjective indicators. Rodrigues et al. (2023) enriched the framework by incorporating social support variables alongside health and frailty measures, using cluster analysis to identify profiles of elderly individuals with distinct care configurations. Their contribution highlighted the importance of the social dimension, but did not extend to digital accessibility or subjective perception.

Han et al. (2025) advanced the multidimensional approach further, proposing a segmentation that integrates health, economic, and social indicators using LCA on European survey data. Their framework represents the closest antecedent to the present study, but retains two significant gaps: the absence of cognitive functioning as a segmentation dimension, and the treatment of subjective well-being measures as outcomes rather than inputs. Across this literature, a consistent pattern emerges: the dimensionality of segmentation has been expanding---from health alone, to health and economics, to health, economics, and social support---but has not yet incorporated the full spectrum of dimensions that characterise elderly vulnerability, and has not treated the subjective experience of ageing as a constitutive element of the segmentation itself.

\section{The Digital Divide in Ageing Populations}
The concept of the digital divide---the unequal distribution of access to, use of, and benefits from digital technologies---has evolved significantly since its initial formulation in the 1990s. In the context of ageing populations, the digital divide acquires particular urgency as public services, healthcare delivery, and social interaction increasingly migrate to digital platforms.

\emph{\textbf{2.2.1 Determinants of digital exclusion among the elderly}}

A substantial body of research has documented the demographic, socioeconomic, and health-related determinants of internet non-use among older adults. Hargittai and Dobransky (2017) demonstrated that age remains the strongest predictor of digital exclusion, even after controlling for education, income, and health status, though the strength of the association varies across institutional contexts. Friemel (2016) found that non-use of the internet among Swiss elderly aged 70 and over was associated with lower education, lower income, poorer health, and living alone---but that the relative importance of these factors differed by gender and cohort. König et al. (2018), using SHARE data, documented the association between internet use and quality of life among European elderly, finding that internet users reported higher life satisfaction even after adjusting for sociodemographic characteristics.

Helsper's (2012) corresponding fields model of digital exclusion provides a useful conceptual bridge: digital inclusion, she argued, is not a property of the individual but of the interaction between individual capabilities (access, skills, attitudes) and the digital demands of specific life fields (health, welfare, social participation). Applied to the Italian case, the framework predicts that digital exclusion should be particularly damaging in fields where public services have been most aggressively digitised---exactly the domain of healthcare after the PNRR Mission 6 reforms. The present study operationalises this prediction by using internet use and related activity indicators as segmentation inputs, rather than treating them as background controls. Anderson and Perrin (2017), reporting Pew Research Center data on technology adoption among older Americans, documented the rapidity with which digital adoption among the over-65 can shift over short cohort windows, a finding that cautions against interpreting today's Italian digital gap as a fixed structural feature rather than a cohort phenomenon with a predictable decay curve.

More recently, research has moved beyond the binary user/non-user distinction to examine the diversity of digital engagement patterns among older adults. Van Deursen and Helsper (2015) proposed a three-level model of digital inequality (access, skills, outcomes), arguing that mere access to the internet is insufficient without the skills to use it effectively and the ability to translate use into tangible benefits. Campens et al. (2024), using two-step cluster analysis on a panel of community-dwelling elderly in the Netherlands, identified distinct typologies of internet use, demonstrating that older adults who go online use the internet for qualitatively different purposes---information seeking, social communication, health management---with different well-being correlates.

\emph{\textbf{2.2.2 Digital exclusion as a health determinant}}

An emerging literature frames digital exclusion not merely as a technology problem but as a social determinant of health. A multi-cohort study across five ageing surveys (including SHARE) by Yuan et al. (2024) found that digital exclusion was significantly associated with higher depressive symptoms among older adults across 24 countries, with the association persisting after adjustment for socioeconomic status, physical health, and social participation. The mechanisms proposed include reduced access to health information, diminished social connectivity, lower engagement with preventive services, and exclusion from digital government services.

For the Italian context, Ferretti et al. (2023) documented the digital health gap among elderly individuals in remote areas of Italy, finding that even basic digital health services---electronic prescriptions, online appointment booking---remained inaccessible for a majority of elderly respondents, particularly those living alone and in peripheral areas. Their findings highlight the interaction between digital exclusion and geographic marginalisation, a dimension that survey data such as SHARE can capture only partially.

Seifert, Cotten and Xie (2021) extended this line of analysis to the COVID-19 pandemic context, finding that the accelerated migration of health services to digital channels widened rather than narrowed pre-existing disparities, with those already excluded falling further behind. This finding is directly germane to the Italian digital health agenda: if the PNRR-financed Fascicolo Sanitario Elettronico follows the pattern Seifert et al. document, its deployment could produce a digital Matthew effect in which those who most need improved healthcare access receive the least benefit from the infrastructure designed to provide it. The profile-based framework developed in this thesis is explicitly designed to identify the segments most exposed to such a dynamic, before policy implementation renders it irreversible.

Tsai, Shillair and Cotten (2015) offered an important methodological caveat. Their review of the internet-use and older-adult well-being literature noted that the selection of healthier and more resourced individuals into internet use renders causal interpretation hazardous, and that much of the literature has treated correlation as though it were effect. This caveat applies directly to the Connected Active profile identified in Chapter 4, whose favourable well-being correlates with but is not necessarily produced by digital engagement. The cross-country design deployed in Chapter 5 is better positioned to address this inferential problem: if within-profile well-being gaps between Italy and Sweden persist for individuals with comparable digital engagement, the residual gap is attributable to institutional rather than individual factors.

Despite this growing body of evidence, the digital divide literature has predominantly treated digital engagement as a dependent variable---something to be predicted or explained---rather than as a dimension on which to segment the population. The present study departs from this convention by integrating digital accessibility as one of five segmentation dimensions, thereby identifying profiles defined partly by their relationship to digital technology rather than merely profiling users versus non-users.

\section{Objective Conditions versus Subjective Well-Being}
A foundational tension in the study of ageing concerns the relationship between objective conditions (health, income, functional capacity) and subjective experience (perceived quality of life, satisfaction, loneliness, hope). The two are correlated but far from identical: an elderly individual with multiple chronic conditions may report high life satisfaction, while a physically healthy individual may experience profound loneliness and despair. Understanding this discrepancy---and its implications for intervention---is essential for any segmentation framework that aspires to capture the full complexity of elderly vulnerability.

\emph{\textbf{2.3.1 The CASP framework}}

The Control, Autonomy, Self-realisation, and Pleasure (CASP) scale, developed by Hyde et al. (2003), represents one of the most theoretically grounded measures of quality of life in older age. Rooted in a needs-satisfaction model derived from Maslow's hierarchy, CASP conceptualises quality of life not as the absence of disease or the presence of material resources, but as the degree to which fundamental human needs---for control over one's environment, autonomy in decision-making, self-realisation through meaningful activity, and the experience of pleasure---are met. The CASP-12 version, a reduced 12-item form validated by Wiggins et al. (2008), has been adopted as the standard quality-of-life measure in SHARE and in several other European longitudinal studies (ELSA, HAPIEE, KLoSA).

Cross-national validation studies have confirmed that the CASP-12 structure is broadly stable across European populations (Oliver et al., 2021), though the autonomy subscale consistently exhibits the lowest reliability. Research using SHARE Wave 7 data (N = 61,355 Europeans aged 60+) has documented that CASP-12 scores decline with age, with the decline accelerating among the oldest old (those aged 85+), and that control and self-realisation show the steepest age gradients while autonomy paradoxically increases slightly in later life (Portellano-Ortiz et al., 2021).

For the purposes of this thesis, the critical observation is methodological: across the SHARE segmentation literature, CASP-12 has been used exclusively as an \emph{outcome} variable---a measure against which to validate clusters or to test hypotheses about quality of life determinants---and never as an \emph{input} to the segmentation itself. The present study breaks with this convention by incorporating CASP-12 as one of the variables on which individuals are clustered, treating perceived quality of life as a constitutive dimension of the profile rather than a consequence to be predicted.

\emph{\textbf{2.3.2 Loneliness: objective isolation versus subjective experience}}

The distinction between social isolation (an objective deficit in social contacts) and loneliness (a subjective experience of inadequate social connection) is well-established in the gerontological literature (Perlman \& Peplau, 1981). The two constructs are moderately correlated but conceptually and empirically distinct: an individual with a large social network may feel lonely if the relationships are perceived as superficial, while a person with few contacts may feel entirely content with their social life. In the SHARE context, social isolation is captured by objective indicators (social network size, social integration index), while loneliness is measured by the UCLA three-item scale (Hughes et al., 2004), a brief validated instrument that asks respondents how often they feel they lack companionship, feel left out, and feel isolated.

The health consequences of loneliness have been extensively documented. Holt-Lunstad et al. (2015), in a meta-analysis of 70 prospective studies (N \textgreater{} 3.4 million), found that social isolation, loneliness, and living alone were each associated with increased mortality risk, with effect sizes comparable to well-established risk factors such as obesity and physical inactivity. Courtin and Knapp (2017), reviewing the evidence on social isolation and loneliness among older Europeans, concluded that both are associated with increased risk of cardiovascular disease, cognitive decline, depression, and functional disability, but that the mechanisms differ: social isolation operates primarily through reduced access to material resources and health-promoting behaviours, while loneliness operates through chronic stress activation and neuroendocrine dysregulation.

A broader reading of the subjective well-being literature reinforces the case for incorporating it into segmentation. Steptoe, Deaton and Stone (2015), in an influential Lancet review, synthesised evidence that subjective well-being is not an epiphenomenon of objective conditions but an independent predictor of longevity, cardiovascular outcomes and healthcare costs, with effect sizes comparable to those of established behavioural risk factors. Treating subjective measures as outcomes of objective conditions, as the prior SHARE segmentation literature has done, therefore misses the reverse causal channel documented in their review.

Diener et al. (2017) extended this argument into a research agenda on measurement, emphasising that subjective well-being is multidimensional---encompassing eudaimonic, hedonic and evaluative components---and that reducing it to a single scale risks masking the heterogeneity of its predictors. The present thesis incorporates this insight by including five subjective variables (CASP, loneliness, hope, interest, survival expectations) that capture distinct facets of this construct, rather than relying on a single life-satisfaction item. Netuveli and Blane (2008) and Bowling (2005) provide the gerontological grounding for this decision. Netuveli and Blane reviewed the measurement of quality of life in older ages, observing that instrument choice has substantive implications: scales anchored in objective functioning produce different rankings from those anchored in lay perspectives or needs-satisfaction models. Bowling's (2005) qualitative work on what elderly respondents themselves understand by 'quality of life' revealed that relationships, autonomy and a sense of purpose consistently outrank functional health in lay accounts---a finding that foregrounds the legitimacy of a segmentation framework in which subjective perception is an input rather than a derived outcome.

A residual methodological concern must be acknowledged. The inclusion of CASP-12 both as a clustering input and as an interpretive dimension of the resulting profiles introduces a degree of circularity: the profiles will, by construction, differ on CASP, so the observed inter-profile CASP gap is partly a mechanical consequence of the design. This thesis addresses the concern on two fronts. The external validation in Chapter 4 relies on life satisfaction and self-reported happiness, neither of which enters the clustering, so the significance tests on those measures are not contaminated by the input-output circularity. And the sensitivity analysis in Chapter 6 compares the five-dimensional solution to a four-dimensional counterfactual that excludes the subjective variables, so the incremental contribution of the subjective dimension can be isolated from its tautological component.

For the present study, the distinction between objective social connectivity and subjective loneliness provides a direct test of the value added by the fifth dimension. If the Factor Analysis reveals that social network size and loneliness load on different factors---that is, if objective social structure and subjective social experience are empirically separable---this constitutes evidence that segmentation based on objective indicators alone misses a dimension of vulnerability that has independent health consequences.

\section{Cognitive Functioning and Digital Engagement}
A less explored but potentially significant relationship concerns the link between cognitive performance and digital engagement in later life. The cognitive reserve hypothesis (Stern, 2002) proposes that education, occupational complexity, and cognitively stimulating activities across the lifespan build a `reserve' that buffers against the clinical expression of neurodegenerative pathology. Individuals with greater cognitive reserve can sustain higher levels of functioning despite equivalent levels of brain pathology, delaying the onset of clinical symptoms.

Several studies have extended this framework to include digital technology use as a form of cognitive stimulation in later life. Xavier et al. (2014) found that internet use among community-dwelling elderly was associated with better performance on standardised cognitive tests, even after controlling for education, age, and socioeconomic status. Opdebeeck et al. (2016), in a meta-analysis, confirmed that engagement in cognitively stimulating activities---a category that plausibly includes internet use, email communication, and information seeking---is associated with reduced risk of cognitive decline. More recently, studies using the Health and Retirement Study (HRS) have documented that social technology use among older adults is associated with lower loneliness, better self-rated health, and fewer depressive symptoms, with the loneliness reduction mediating a substantial portion of the health benefits (Chopik, 2016).

The evidence base on digital engagement as cognitive stimulation is, however, contested. Small et al. (2009) demonstrated, using functional MRI, that internet searches activate brain regions associated with decision-making and complex reasoning in novices, with patterns approaching those of regular users after only a few days of practice; the authors interpreted this as evidence that the cognitive demands of internet use are themselves non-trivial. Klusmann et al. (2010), in one of the few randomised controlled trials of computer-based cognitive training in older adults, found statistically significant though modest improvements in episodic memory and executive function among participants assigned to a six-month training programme, relative to controls. Czaja et al. (2006), in the foundational CREATE consortium studies, documented that older adults do not lack capability but face instruments designed with assumptions that fail to accommodate their needs---reframing the cognitive-digital question as one of interaction design rather than of inherent ability.

A more sceptical voice must be registered. Boot et al. (2015) reviewed the claims made for 'brain training' and related cognitive-enhancement products, concluding that the published evidence is vulnerable to placebo effects, publication bias, and the near-universal finding that transfer of training to untrained cognitive domains is weak. The relevance of this critique to the present thesis is direct: the digital-cognitive co-loading identified in the Factor Analysis of Chapter 4 should not be interpreted as evidence that internet use causally improves cognition. It indicates only that the two capacities covary in later life, which is a substantially weaker claim. The cross-country comparison of Chapter 5 supports this cautious reading: Swedish elderly outperform Italians on verbal fluency by a magnitude that cannot plausibly be attributed to internet-use differentials alone, and the gap persists in matched profiles where digital engagement is comparable.

The question of whether cognitive performance and digital engagement reflect a single underlying construct---a shared `cognitive-digital capital'---or two correlated but distinct dimensions has not been directly addressed in the segmentation literature. The present study contributes to this question through the Factor Analysis, which examines whether cognitive variables (verbal fluency) and digital variables (internet use) load on the same or different latent factors. A co-loading would suggest that interventions targeting cognitive stimulation and digital inclusion may be addressing the same underlying capacity, with implications for programme design and resource allocation.

\section{Welfare Regimes and Cross-Country Comparisons}
The comparative welfare state literature provides the theoretical framework for the cross-country component of this thesis. Esping-Andersen's (1990) typology of welfare regimes---liberal, conservative-corporatist, and social-democratic---remains the foundational reference, later enriched by Ferrera's (1996) identification of the Southern European model as a distinct regime type characterised by fragmented social insurance, minimal public services, and heavy reliance on family solidarity. The implications for the elderly are direct: in the social-democratic model (Sweden), the state provides extensive home help, senior activity centres, and institutional support for independent living; in the Southern European model (Italy), these functions are primarily delegated to the family, with public provision serving a residual role.

Diener, E., Heintzelman, S. J., Kushlev, K., Tay, L., Wirtz, D., Lutes, L. D., \& Oishi, S. (2017). Findings all psychologists should know from the new science on subjective well-being. Canadian Psychology, 58(2), 87--104.

The cross-country comparison of elderly well-being has been pursued using SHARE data by several authors, though typically at the aggregate level or with a focus on specific outcomes. Börsch-Supan et al. (2013), in the original SHARE data resource profile, documented systematic cross-national variation in health, cognition, and economic resources among Europeans aged 50+. Subsequent studies have compared specific dimensions---depression prevalence (Castro-Costa et al., 2008), cognitive performance (Skirbekk et al., 2012), or healthcare utilisation (Bonsang, 2009)---across countries, generally finding that Northern European countries exhibit better outcomes on most dimensions, with the notable exception of wealth inequality.

Esping-Andersen's tripartite typology has been productively contested and refined. Bambra (2007) noted that the original schema, constructed on the basis of pension and unemployment schemes, generalises poorly to care-related social services, where cross-country variation does not align cleanly with the three-regime model; empirical clustering of care indicators tends to produce a four- or five-regime solution in which Italy and Sweden remain poles apart even as intermediate cases redistribute. Bonoli (1997) proposed a two-dimensional refinement distinguishing the 'how much' (social expenditure) from the 'how' (Bismarckian versus Beveridgean principles of entitlement), a framework that helps explain why Italian healthcare can be formally universalist while producing the cost-driven access gradients documented in Chapter 7 of this thesis. Anttonen and Sipilä (1996), in an early and influential analysis specifically focused on social care services, identified Scandinavian countries as the most service-intensive and Mediterranean countries as the most family-intensive, a contrast that maps directly onto the Moderate Isolated--Traditional Social divergence across the two countries examined here. Bettio and Plantenga (2004) later documented that migrant-worker-based care arrangements---the so-called badanti phenomenon in Italy---represent a market-based substitute for the public care services characteristic of the Scandinavian model, without providing the equivalent of the information-circulation and social-contact functions that Chapter 7 identifies as the principal channel through which social connectivity mediates healthcare access.

The empirical magnitude of the Italy--Sweden divergence at the welfare-regime level is stark. Public expenditure on long-term care as a share of GDP is approximately 3.2\% in Sweden and 0.9\% in Italy (OECD, 2023), a threefold gap that persists after accounting for differences in age structure and for the monetary valuation of informal care. The coverage of publicly funded home help among those aged 80 and over is approximately 25\% in Sweden and below 5\% in Italy. Swedish municipalities provide universal access to senior activity centres, mediated through municipality-level social-services offices that double as information points for entitlements; the Italian equivalent, where it exists at all, is largely delegated to the third sector and voluntary associations, whose geographic coverage is uneven and whose clientele is systematically drawn from the already-connected. This institutional asymmetry is the system-level reason the Moderate Isolated profile is a 28.5\% segment of the Italian elderly population and has no structural counterpart in Sweden. Contrary to the aggregate-level framings that dominate the comparative welfare literature, the present thesis establishes this result at profile level: it is not only that Italy and Sweden differ on average, but that the same types of elderly experience radically different material and subjective conditions depending on the welfare regime in which they age.

What has not been attempted, to the best of this author's knowledge, is a multidimensional \emph{profile-level} comparison between a Southern and a Northern European country using a common five-dimensional segmentation framework that includes subjective perception variables. The existing cross-country SHARE literature compares aggregate means or single-dimension profiles, but does not ask the question: do the same types of elderly individuals exist in both countries, and if so, do they experience their condition differently? This question---which lies at the intersection of population segmentation and comparative welfare state analysis---is the central comparative contribution of this thesis.

\section{The Expectations Module: An Unexplored Resource}
A distinctive feature of the present study is the inclusion of variables from SHARE's Expectations module (EX), which has received minimal attention in the segmentation literature. The EX module contains questions about subjective probabilities of future events: the probability of surviving to a target age (ex001\_), the probability that the government will reduce pensions (ex007\_), and the probability of receiving an inheritance (ex009\_). These variables capture an individual's orientation toward the future---their expectations, hopes, and perceived risks---which are conceptually distinct from both current conditions and current feelings.

The survival probability question (ex001\_) is particularly interesting: an elderly person who assigns a 20\% probability to surviving another decade is operating under a fundamentally different horizon than one who assigns 80\%. This orientation plausibly affects health behaviours, social investment, financial decisions, and engagement with services, yet it has not been incorporated into any SHARE segmentation study. The present thesis includes ex001\_ (recoded as expect\_alive) as one of the five subjective perception variables, testing whether survival expectations contribute meaningfully to the differentiation of elderly profiles.

\section{Research Gaps and Positioning of the Present Study}
The review of the literature reveals five interconnected gaps that the present study addresses.

\begin{table}[ht]
  \centering
  \caption{Research gaps and contributions of the present study.}
  \label{tab:ch2-gaps}
  \small
  \begin{tabular}{@{}p{3.5cm}p{5cm}p{5cm}@{}}
    \toprule
    \textbf{Gap} & \textbf{Existing literature} & \textbf{This thesis}\\
    \midrule
    Subjective variables as segmentation inputs & CASP, loneliness, expectations used as outcomes only & CASP, loneliness, hope, interest, expect\_alive as clustering inputs\\
    Five-dimensional framework & 2--3 dimensions (health, economics, social) & 5 dimensions: health, economics, digital/social, cognitive, subjective\\
    Cross-country profile comparison & Aggregate means or single-dimension profiles & Profile-level comparison with matched segments (IT vs SE)\\
    Digital dimension as segmentation input & Digital use as dependent variable to be predicted & Digital accessibility as one of five clustering dimensions\\
    Expectations module in segmentation & EX module not used in any segmentation study & \emph{expect\_alive} included as subjective perception variable\\
    \bottomrule
  \end{tabular}
  \par\smallskip
  \footnotesize \emph{Note.} IT = Italy, SE = Sweden. See Chapter~\ref{ch:ch3} for the complete variable specification.
\end{table}

First, the SHARE segmentation literature has employed subjective well-being measures exclusively as validation outcomes, never as segmentation inputs. This thesis incorporates CASP-12, loneliness, hope for the future, interest in things, and survival expectations as variables on which individuals are clustered, treating the subjective experience of ageing as a constitutive dimension of the profile rather than a consequence to be explained.

Second, the existing segmentation frameworks have employed at most three dimensions (health, economics, social support). This thesis extends the framework to five dimensions by adding cognitive functioning and subjective perception as distinct analytical axes, enabling the identification of profiles that are invisible in lower-dimensional representations---most notably, the distinction between two types of frailty (resigned and depressed) that are indistinguishable on objective indicators alone.

Third, cross-country comparisons on SHARE have focused on aggregate indicators or single-dimension analyses. This thesis conducts a profile-level comparison, asking not `do Italians and Swedes differ on average?' but `do the same types of elderly exist in both countries, and if so, do matched profiles exhibit different levels of subjective well-being and social connectivity?' This profile-matching approach provides a more granular and actionable comparison than aggregate statistics.

Fourth, the digital divide literature has treated digital engagement as a dependent variable. This thesis integrates it as a segmentation dimension, enabling the identification of profiles defined partly by their digital accessibility---a directly relevant framing for the policy question of how the digitalisation of public services interacts with the heterogeneity of the elderly population.

Fifth, the SHARE Expectations module has not been used in any segmentation study. By including survival expectations as one of the subjective perception variables, this thesis tests whether an individual's orientation toward the future contributes meaningfully to profile differentiation---a theoretically motivated but empirically untested proposition.

Additionally, from a methodological standpoint, the use of Latent Class Analysis alongside PCA-based k-means clustering and Factor Analysis provides a degree of methodological triangulation that is not standard in the course curriculum and strengthens confidence in the identified profiles through the convergence of fundamentally different analytical paradigms. Taken together, these five gaps are not incremental refinements of the existing literature; they reflect a cumulative choice about what a segmentation of the elderly is for. The literature surveyed in this chapter has progressively expanded the dimensional reach of such segmentations, from health alone, to health and economics, to health, economics and social support. The step taken here---adding subjective perception as an input rather than an outcome, and validating the entire framework across two welfare regimes---completes that trajectory rather than departing from it.
